\chapter{Biological Background}
\label{chap:background}

This chapter provides the biological foundation required to understand the
problem that PROTEA addresses. Readers with a background in molecular biology
may skip directly to \cref{chap:related}.

% ────────────────────────────────────────────────────────────
\section{Proteins: Structure and Function}
\label{sec:proteins}

Proteins are macromolecules composed of chains of amino acids linked by peptide
bonds. The linear sequence of amino acids — encoded by the genome — is referred
to as the \textit{primary structure}. This sequence folds, driven by
thermodynamic forces, into a specific three-dimensional conformation that
determines the protein's biological activity.

Protein functions are extraordinarily diverse: enzymes catalyse chemical
reactions, structural proteins provide mechanical support, signalling proteins
relay information between cells, and transport proteins shuttle molecules across
membranes. Understanding \textit{what} a protein does — its molecular function,
the biological process it participates in, and the cellular compartment where it
operates — is fundamental to biology and medicine.

\subsection{Amino Acids and the Sequence Space}
\label{subsec:amino-acids}

Twenty standard amino acids, differing in their side-chain chemistry, constitute
the alphabet of protein sequences. A protein of $n$ residues can in principle
take $20^n$ distinct sequences, an astronomically large space. Yet the sequences
that fold into stable, functional structures represent a tiny, structured subset
of this space, one in which similar sequences tend to adopt similar folds and
perform similar functions.

\subsection{Sequence Identity and Homology}
\label{subsec:homology}

Two proteins are said to be \textit{homologous} if they share a common
evolutionary ancestor. Homology is typically inferred from sequence similarity:
proteins with high pairwise sequence identity are almost always homologous and
usually perform the same function. As sequence identity drops below roughly
30\,\%, however, the relationship between similarity and function becomes
ambiguous — proteins may be structurally similar yet functionally diverged, or
functionally convergent with unrelated folds.

% ────────────────────────────────────────────────────────────
\section{Protein Databases}
\label{sec:databases}

\subsection{UniProt}
\label{subsec:uniprot}

\gls{uniprot} is the central repository for protein sequence and functional
information~\cite{uniprot2023}. It is divided into two sections:

\begin{description}
  \item[\gls{swissprot}] A manually curated, high-quality database. Each entry
    has been reviewed by expert annotators and contains experimental or
    computationally inferred functional annotations, with explicit evidence codes.
  \item[\gls{trembl}] An automatically annotated, computationally analysed
    database derived from genome translations. It is several orders of magnitude
    larger than Swiss-Prot but of lower average annotation quality.
\end{description}

Each protein in UniProt is identified by a stable accession number. Canonical
isoforms are represented by a bare accession (e.g.\ \texttt{P12345}); alternative
splicing variants are denoted \texttt{P12345-2}, \texttt{P12345-3}, and so on.

\subsection{Protein Data Bank}
\label{subsec:pdb}

The \gls{pdb} archives three-dimensional structures of biological macromolecules
determined experimentally by X-ray crystallography, cryo-electron microscopy, or
NMR~\cite{berman2000pdb}. Although PROTEA operates primarily on sequences, PDB
structures are an important source of ground-truth functional information and are
used indirectly through UniProt cross-references.

% ────────────────────────────────────────────────────────────
\section{The Gene Ontology}
\label{sec:gene-ontology}

\subsection{Overview}
\label{subsec:go-overview}

The \gls{go} is a standardised, species-independent controlled vocabulary for
describing gene product attributes~\cite{ashburner2000go,goconsortium2023}.
It is structured as three independent \glspl{dag}, one for each of the three
top-level \textit{aspects}:

\begin{description}
  \item[\gls{mf} (GO:0003674)] The biochemical activity of the gene product,
    independent of where or when it occurs (e.g.\ \textit{ATP binding},
    \textit{serine-type endopeptidase activity}).
  \item[\gls{bp} (GO:0008150)] The larger biological programme to which the
    activity contributes (e.g.\ \textit{apoptotic process},
    \textit{DNA repair}).
  \item[\gls{cc} (GO:0005575)] The place in the cell where the gene product
    is active (e.g.\ \textit{nucleus}, \textit{plasma membrane}).
\end{description}

Terms within each aspect are connected by \textit{is-a} and \textit{part-of}
relationships, forming a hierarchy from broad root terms to highly specific
leaves. A protein annotated with a term is implicitly annotated with all of its
ancestors up to the root — the \textit{true path rule}.

\subsection{Evidence Codes}
\label{subsec:evidence-codes}

Every GO annotation carries an evidence code that indicates the type of support
for the association. Codes range from experimental (e.g.\ \texttt{EXP},
\texttt{IDA}, \texttt{IMP}) to computationally inferred (\texttt{ISS}, \texttt{IEA})
and author-stated (\texttt{TAS}, \texttt{NAS}). Experimental evidence codes are
the gold standard; \texttt{IEA} (Inferred from Electronic Annotation) is the most
common but least reliable.

\subsection{Ontology Versioning}
\label{subsec:go-versioning}

The GO is actively maintained and released regularly in \gls{obo} format.
New terms are added, obsolete terms are deprecated, and relationships are
refined with each release. Reproducible research therefore requires that
analyses record which ontology version was used. PROTEA stores each imported
release as an \texttt{OntologySnapshot} and associates all annotation sets
and predictions with a specific snapshot.

% ────────────────────────────────────────────────────────────
\section{Sequence Alignment}
\label{sec:alignment}

Pairwise sequence alignment is the classical approach to measuring sequence
similarity. Two algorithms are central to this thesis:

\begin{description}
  \item[Needleman--Wunsch (\gls{nw})] A global alignment algorithm that aligns
    two sequences end-to-end, maximising a cumulative substitution score minus
    gap penalties~\cite{needleman1970general}. It is suitable for comparing
    proteins of similar length.
  \item[Smith--Waterman (\gls{sw})] A local alignment algorithm that finds the
    highest-scoring contiguous aligned region, ignoring the tails of each
    sequence~\cite{smith1981identification}. It is more appropriate when one
    sequence is a domain or fragment of the other.
\end{description}

Both algorithms use a \textit{substitution matrix} — typically BLOSUM62 for
protein sequences~\cite{henikoff1992amino} — to score residue substitutions
based on their observed frequency in alignments of related proteins.

% ────────────────────────────────────────────────────────────
\section{Protein Language Models and Embeddings}
\label{sec:plm}

\Glspl{plm} are neural networks trained on large corpora of protein sequences
using self-supervised objectives borrowed from natural language processing.
Models such as ESM-2~\cite{lin2023esm2} and ProtTrans~\cite{elnaggar2022prottrans}
learn to predict masked residues, in doing so developing internal representations
that capture evolutionary, structural, and functional information.

A protein sequence can be passed through such a model to obtain a dense vector
(embedding) that encodes its properties in a continuous latent space. Proteins
with similar functions tend to cluster together in this space even when their
sequences have diverged below the homology detection threshold of alignment-based
tools. This property makes embeddings attractive for annotation transfer.

% ────────────────────────────────────────────────────────────
\section{Taxonomic Classification}
\label{sec:taxonomy}

Biological organisms are classified into a hierarchy — the NCBI taxonomy — that
reflects their evolutionary relationships~\cite{schoch2020ncbi}. Each node is
assigned a numeric \textit{taxon ID}. The distance between two organisms can be
defined operationally as the number of edges between them through their
\gls{lca} in the taxonomy tree.

Taxonomic distance is a useful co-variate for function prediction: a protein from
a closely related organism is more likely to perform the same function than one
from a distant lineage, even when embedding similarity is comparable.
